\documentclass[12pt]{report}
\usepackage{float}
\usepackage[a4paper, margin=0.9in]{geometry}
\usepackage{graphicx}
\usepackage{hyperref}
\usepackage{wrapfig}
\usepackage{acronym}
\usepackage{rotating}
\usepackage{glossaries}
\usepackage{mdframed}
\usepackage{enumitem}
\usepackage[table]{xcolor}
\usepackage[raster]{tcolorbox}
\usepackage{caption}
\usepackage[french]{babel}
\usepackage{enumitem} % Pour personnaliser les listes
\usepackage{setspace}  % Added for line spacing control
\usepackage[T1]{fontenc}
\usepackage{amsmath,amssymb}
\definecolor{swotS}{RGB}{226,237,143}
\definecolor{swotW}{RGB}{247,193,139}
\definecolor{swotO}{RGB}{173,208,187}
\definecolor{swotT}{RGB}{192,165,184}
\usepackage{listings}
\usepackage{xcolor}
\lstset{
    basicstyle=\ttfamily\small,
    keywordstyle=\color{blue},
    commentstyle=\color{orange},
    stringstyle=\color{red},
    breaklines=true,
    numbers=left,
    numberstyle=\tiny\color{gray},
    frame=single,
    captionpos=b
}
\hypersetup{
    colorlinks=true,
    linkcolor=blue,
    filecolor=magenta,
    urlcolor=cyan,
}

\begin{document}
\onehalfspacing  % Adjust line spacing for the document

\begin{titlepage}

\begin{figure}[h]
    \centering
\begin{center}
\begin{minipage}{0.2\textwidth}
    \centering
    \includegraphics[width=1.3\linewidth]{ENSABERRECHID_logo20230124170433.png}
\end{minipage}\hfill
\begin{minipage}{0.58\textwidth}
    \centering
    {\footnotesize \bf ROYAUME DU MAROC\\
    UNIVERSITÉ HASSAN PREMIER\\
    ÉCOLE NATIONALE DES SCIENCES APPLIQUÉES BERRECHID}
\end{minipage}\hfill
\begin{minipage}{0.1\textwidth}
    \centering
    \includegraphics[width=1.3\linewidth]{pfe/fst-hassan-1er (1).png}
\end{minipage}
\end{center}
\end{figure}

\begin{center}
\vspace{2cm}
%------------------------------------------------------------------------- 06
\Huge{\textcolor{red}{\textbf{Rapport de Projet}}}\\
\vspace{0.5cm}
\Large{\textit{Spécialité:} \textbf{ Ingénierie des Systèmes d'Information et Big Data}\\ 
\Large{\textit{Module:} \textbf{ Systéme Business Intelligence}\\ 
\noindent\rule{\textwidth}{1mm}
\HUGE{\textbf{Maintenance prédictive pour IoT industriel avec Data Warehouse temporel}}
\noindent\rule{\textwidth}{1mm}
\end{center}

\begin{center}
\end{center}

%-------------------------------------------------------------------------- 07
\begin{center}
\begin{tabular}{ p{11.5cm} p{15cm} }
\Large{\textcolor{red}{\textit{Réalisé par:}}} & \Large{\textcolor{red}{\textit{Encadré par:}}} \\
\vspace{0.01cm} \textbf{FARMATI Aya } & \vspace{0.03cm} \textbf{Prof. HRIMECH Hamid} \\
\vspace{0.001cm} \textbf{KHACHANE Zakaria }
\end{tabular}
\end{center}

\vspace{1cm}

\begin{center}
\\
\end{center}
\begin{center}
\vspace{0.3cm}

%--------------------------------------------------------------------------- 08
\vspace{3.5cm}

\begin{center}
\Large{\textit{Année Universitaire: 2025-2026}}
\end{center}

\end{titlepage}

\begin{center}
\addcontentsline{toc}{chapter}{Remerciement}
\section*{Remerciement}
\end{center}

\begin{center}
  \par\nobreak\vspace{\baselineskip}
  \noindent\rule{12cm}{0.2pt}
  \par\nobreak\vspace{\baselineskip}
\end{center}

{\Large\itshape
Avant toute chose, il apparaît opportun de commencer ce rapport par des remerciements sincères, à l’adresse de tous ceux qui nous ont soutenus, encouragés et transmis leur savoir tout au long de notre parcours universitaire.

\par\bigskip

Nous tenons d’abord à remercier Dieu le Tout-Puissant, qui nous a donné la force, la persévérance et la patience nécessaires pour mener à bien ce travail et surmonter les difficultés rencontrées.

\par\bigskip

Nous exprimons également toute notre gratitude à l’ensemble de nos enseignants, pour la qualité de leur enseignement et leur engagement constant à nos côtés. Plus particulièrement, nous adressons nos remerciements chaleureux à notre encadrant, le professeur HRIMECH Hamid, pour ses conseils avisés, sa disponibilité et ses encouragements précieux, qui ont largement contribué à la réussite de ce projet.

\par\bigskip

Enfin, nous remercions profondément nos parents et nos proches pour leur soutien moral indéfectible et leurs encouragements continus, qui ont été pour nous une source inépuisable de motivation et de réconfort.

\par\bigskip
}

\begin{center}
  \par\nobreak\vspace{\baselineskip}
  \noindent\rule{12cm}{0.2pt}
  \par\nobreak\vspace{\baselineskip}
\end{center}

\tableofcontents
\addcontentsline{toc}{chapter}{Table des figures}
\listoffigures

\begin{center}
\addcontentsline{toc}{chapter}{Introduction Générale}
\chapter*{Introduction Générale}
\end{center}

Dans le cadre de ma formation en Ingénierie des Systèmes d'Information et Big Data, ce projet de fin de module en Business Intelligence représente une opportunité clé pour mettre en œuvre des architectures décisionnelles modernes. Le domaine de la maintenance prédictive industrielle est aujourd'hui en pleine expansion, propulsé par la prolifération des capteurs connectés (IoT) et les technologies d'analyse de données massives (Big Data).

L'objectif de ce projet est de concevoir et développer une plateforme complète (End-to-End) de collecte, de stockage, d'analyse prédictive et de visualisation décisionnelle pour la maintenance prédictive de moteurs d'avion (turbofans) à partir de données simulées. La finalité est d'estimer en temps réel la Durée de Vie Utile Restante (\textit{RUL - Remaining Useful Life}) afin d'éviter les pannes critiques et d'optimiser les cycles de maintenance.

Ce système s'appuie sur des technologies décisionnelles et MLOps de pointe :
\begin{itemize}
    \item Une ingestion temps réel de flux IoT via le protocole \textbf{MQTT} et \textbf{Apache Spark Structured Streaming}.
    \item Un stockage Lakehouse distribué sur \textbf{MinIO} structuré en couches \textbf{Medallion (Bronze, Silver, Gold)} sous le format transactionnel \textbf{Apache Hudi}.
    \item Un cycle de vie de modélisation prédictive basé sur \textbf{XGBoost} et tracé par le registre de modèles \textbf{MLflow}.
    \item Un Data Warehouse dimensionnel en schéma en étoile hébergé sur \textbf{ClickHouse} pour des analyses de Business Intelligence (BI) ultra-rapides.
    \item Une orchestration centralisée et robuste de l'ensemble des flux de calcul gérée par \textbf{Apache Airflow} en mode événementiel.
\end{itemize}

Ce rapport vise à documenter de manière exhaustive le travail de conception et d'implémentation réalisé à travers cinq chapitres principaux :
\begin{itemize}
\item \textbf{Chapitre 1 :} Présentation de l'architecture générale du projet et de la stack technique.
\item \textbf{Chapitre 2 :} Modélisation prédictive MLOps avec XGBoost et MLflow.
\item \textbf{Chapitre 3 :} Modélisation dimensionnelle (Gold) et intégration avec Power BI.
\item \textbf{Chapitre 4 :} Détection de dérives statistiques (Concept \& Data Drift) en temps réel.
\item \textbf{Chapitre 5 :} Orchestration globale des pipelines de données, CI/CD et déploiement.
\end{itemize}


\chapter{Architecture Générale du Projet}

\section{Modélisation du Flux de Données Medallion}
Le projet implémente une architecture Lakehouse de type Medallion, découpée en trois couches logiques distinctes garantissant la qualité et la traçabilité de la donnée :
\begin{itemize}
    \item \textbf{Couche Bronze (Ingestion brute)} : Reçoit les messages JSON des capteurs de température, pression, vitesse physique et opérationnelle des moteurs par le protocole MQTT. Les données sont insérées en continu dans MinIO au format Apache Hudi en mode \textit{Merge-On-Read (MOR)} afin de supporter des écritures haute fréquence à faible latence.
    \item \textbf{Couche Silver (Feature Store)} : Traite en batch les données brutes. Ce processus applique un pivotement pour transformer le flux (au format clé-valeur) en schéma relationnel large, lisse le bruit des capteurs via des agrégations glissantes par minute, et impute les éventuelles valeurs manquantes à l'aide d'un référentiel historique de référence (\textit{sensor\_baselines}).
    \item \textbf{Couche Gold (DW Dimensionnel)} : Exécute l'inférence prédictive sur les caractéristiques nettoyées et structure les résultats en schéma décisionnel en étoile pour l'alimentation de tableaux de bord de Business Intelligence.
\end{itemize}

\section{Stack Technique}
L'ensemble de l'écosystème s'appuie sur des outils phares du Big Data et du décisionnel :
\begin{itemize}
    \item \textbf{MQTT (Eclipse Mosquitto)} : Broker de messagerie léger pour la réception des flux de capteurs.
    \item \textbf{Apache Spark} : Utilisé pour le streaming continu (Bronze) et le traitement distribué batch (Silver $\rightarrow$ Gold).
    \item \textbf{Apache Hudi} : Format de stockage transactionnel (ACID, Schema Evolution) au-dessus de MinIO.
    \item \textbf{ClickHouse} : Base de données relationnelle colonnaire (OLAP) optimisée pour le requêtage analytique temps réel.
    \item \textbf{MLflow} : Outil de cycle de vie ML pour la traçabilité des runs et le registre des modèles.
    \item \textbf{Apache Airflow} : Système d'orchestration pour l'automatisation et le monitoring des workflows.
\end{itemize}


\chapter{Modélisation Prédictive et MLOps}

\section{Préparation des Données et Extraction de Features}
Le modèle prédictif s'appuie sur l'historique temporel immédiat des capteurs pour prédire le RUL. Des indicateurs statistiques glissants de moyenne et d'écart-type sont générés sur des fenêtres de \textbf{5, 10 et 15 cycles} pour capturer les tendances d'usure des pièces mécaniques :
\begin{equation}
\mu_{\text{capteur}, w} = \frac{1}{w}\sum_{i=0}^{w-1} x_{t-i}
\end{equation}
\begin{equation}
\sigma_{\text{capteur}, w} = \sqrt{\frac{1}{w}\sum_{i=0}^{w-1} (x_{t-i} - \mu_{\text{capteur}, w})^2}
\end{equation}

Pour éviter que le modèle ne tente d'apprendre des tendances de dégradation inexistantes au début de la vie du moteur (où l'usure est nulle et le bruit domine), la cible de RUL est plafonnée à \textbf{130 cycles} dans le script d'entraînement via la formule :
\begin{equation}
y = \min(\text{RUL}_{\text{réel}}, 130)
\end{equation}

\section{Entraînement et Tracking sous MLflow}
Un modèle de régression XGBoost (\texttt{XGBRegressor}) est entraîné avec les hyperparamètres suivants :
\begin{itemize}
    \item Nombre d'estimateurs (\texttt{n\_estimators}) : 500
    \item Profondeur maximale (\texttt{max\_depth}) : 6
    \item Taux d'apprentissage (\texttt{learning\_rate}) : 0.03
    \item Sous-échantillonnage (\texttt{subsample}) : 0.8
\end{itemize}

Le script \texttt{train\_xgboost.py} enregistre l'ensemble des runs sur MLflow. Les métriques de performance obtenues (RMSE, MAE) sont tracées, la signature des entrées/sorties du modèle est définie et le modèle est enregistré sous le nom \texttt{rul\_xgboost}. L'échantillon des données d'entraînement (\texttt{X\_train}) est exporté au format Parquet sur MinIO en tant que référence statistique.


\chapter{Modélisation Dimensionnelle et Power BI}

\section{Schéma en Étoile ClickHouse}
Pour permettre des analyses Business Intelligence instantanées, la couche Gold structure les données en un schéma dimensionnel en étoile implémenté dans ClickHouse :

\begin{figure}[H]
    \centering
    \begin{tikzpicture}[node distance=2.2cm, every node/.style={fill=blue!10, draw, rounded corners, font=\footnotesize, align=center}]
        % Central Fact Table
        \node[fill=red!20, minimum width=4cm, minimum height=2cm] (fact) {\textbf{fact\_engine\_health} \\ \hline fact\_id (PK) \\ engine\_id (FK) \\ date\_id (FK) \\ status\_id (FK) \\ time\_in\_cycles \\ predicted\_rul \\ sensor\_metrics...};
        
        % Dimensions
        \node[left=of fact, minimum width=3cm] (dim_engine) {\textbf{dim\_engine} \\ \hline engine\_id (PK) \\ sensor\_id \\ unit\_number \\ engine\_model};
        \node[above=of fact, minimum width=3cm] (dim_date) {\textbf{dim\_date} \\ \hline date\_id (PK) \\ timestamp \\ year, month, day \\ hour, minute};
        \node[right=of fact, minimum width=3cm] (dim_status) {\textbf{dim\_status} \\ \hline status\_id (PK) \\ status\_label \\ description};
        
        % Relations
        \draw[->, thick] (fact) -- (dim_engine);
        \draw[->, thick] (fact) -- (dim_date);
        \draw[->, thick] (fact) -- (dim_status);
    \end{tikzpicture}
    \caption{Schéma en étoile de la couche Gold}
\end{figure}

Les tables de dimension sont définies de la manière suivante :
\begin{itemize}
    \item \textbf{dim\_engine} : Contient les informations relatives aux moteurs physiques.
    \item \textbf{dim\_date} : Contient le découpage temporel jusqu'à la minute près.
    \item \textbf{dim\_status} : Permet de décoder le niveau d'alerte calculé d'après le RUL :
    \begin{itemize}
        \item \textit{Critique} (RUL $\le$ 15) : Maintenance immédiate.
        \item \textit{Avertissement} (15 < RUL $\le$ 45) : Planifier une maintenance.
        \item \textit{Sain} (RUL > 45) : Fonctionnement normal.
    \end{itemize}
\end{itemize}

\section{Connexion et Modélisation dans Power BI}
Le Data Warehouse ClickHouse est directement connecté à \textbf{Power BI Desktop} via le pilote \textbf{ClickHouse ODBC (Unicode)} ou le connecteur ClickHouse natif. La configuration de la connexion de données s'appuie sur la chaîne de connexion suivante :
\begin{lstlisting}[caption=Configuration de la chaîne de connexion ODBC]
Driver={ClickHouse ODBC (Unicode)};Server=localhost;Port=8123;Database=iot_metrics_DW;Uid=VOTRE_USER;Pwd=VOTRE_MDP;
\end{lstlisting}

Dans Power BI, les tables sont importées (\textit{Mode Import}) ou interrogées directement (\textit{Mode DirectQuery}) pour un rafraîchissement temps réel lors de l'ingestion MQTT. Les relations entre la table des faits \texttt{fact\_engine\_health} et les dimensions \texttt{dim\_engine}, \texttt{dim\_date} et \texttt{dim\_status} sont matérialisées dans l'onglet de modélisation de Power BI (relations un-à-plusieurs directionnelles). Le fichier modèle est fourni directement dans le dépôt sous le dossier \texttt{PowerBi/} sous le nom de \texttt{Analyse\_Maintenance\_Predictive\_Moteurs.pbix}.

\section{Rapports et Visualisations BI Construits}
Le tableau de bord interactif \textbf{"Analyse Temporelle de la Maintenance Prédictive des Moteurs"} a été structuré en cinq pages d'analyse spécialisées :

\begin{enumerate}
    \item \textbf{Overview (Vue d'ensemble)} :
    \begin{itemize}
        \item Filtres interactifs et slicers de navigation par \textit{Moteur}, \textit{Date} et \textit{Etat du Moteur}.
        \item Graphique linéaire représentant le \textit{RUL Prédit dans le Temps}.
        \item Diagramme circulaire (\textit{Pie Chart}) illustrant la répartition des \textit{Etats des Moteurs Actuels} (Sain, Avertissement, Critique).
        \item Tableau synthétique \textit{Résumé des Metrics Actuellement} présentant les dernières valeurs de capteurs du réacteur sélectionné.
    \end{itemize}
    
    \item \textbf{Metrics dans le Temps (Pages 1 \& 2)} :
    Permet de suivre de manière dynamique les courbes temporelles de signaux physiques pour corréler l'usure mécanique avec les estimations d'espérance de vie du modèle :
    \begin{itemize}
        \item \textit{Page 1} : Graphiques temporels pour le RUL prédit, la température de sortie du compresseur HP, la température de sortie du compresseur BP, la pression totale du compresseur HP, la température de sortie de la turbine BP, et le ratio carburant/pression.
        \item \textit{Page 2} : Suivi de la pression statique du compresseur HP, des vitesses physique et corrigée du cœur, de l'enthalpie d'air de purge, et des débits de refroidissement des turbines HP et BP.
    \end{itemize}
    
    \item \textbf{Analyse des Corrélations (Pages 1/2 \& 2/2)} :
    Permet de mettre en évidence la relation statistique directe entre les mesures physiques des capteurs (axe X) et le RUL prédit (axe Y) à l'aide de nuages de points (\textit{Scatter Plots}) munis de courbes de tendance linéaire (\textit{Trendlines}) :
    \begin{itemize}
        \item \textit{Correlation 1/2} : Analyse des variables et capteurs de la page de métriques 1.
        \item \textit{Correlation 2/2} : Analyse des variables et capteurs de la page de métriques 2.
        \item \textit{Interprétation physique} : Les températures et pressions de compresseurs présentent une pente descendante (corrélation négative avec le RUL, signifiant une hausse de friction thermique à l'approche de la panne), tandis que les débits de refroidissement des turbines ont une pente ascendante (corrélation positive avec le RUL).
    \end{itemize}
\end{enumerate}


\chapter{Surveillance de Dérive de Données (Drift)}

\section{Fondement Statistique}
Le modèle perd en pertinence si la distribution des données réelles dévie de celle des données d'entraînement. Nous implémentons un algorithme de détection de \textbf{Data Drift} basé sur le test non-paramétrique bidimensionnel de \textbf{Kolmogorov-Smirnov (KS-test)}.

Pour chaque variable de capteur, la statistique de test $D$ est calculée par :
\begin{equation}
D = \sup_{x} |F_{\text{ref}}(x) - F_{\text{courant}}(x)|
\end{equation}
Où $F_{\text{ref}}$ et $F_{\text{courant}}$ sont les fonctions de répartition empiriques respectives des données d'entraînement (référence) et des données de production actuelles (Silver).

\section{Détection du Concept Drift}
Une variable est déclarée comme dérivée si la p-value associée à la statistique $D$ est inférieure à un seuil critique $\alpha = 0.05$. Le système évalue le pourcentage global de caractéristiques ayant dérivé :
\begin{equation}
\text{Part Drift} = \frac{\text{Nombre de variables dérivées}}{\text{Nombre de variables totales}}
\end{equation}

Si cette part dépasse 20\% ($\text{Part Drift} \ge 0.20$), le drapeau global \texttt{dataset\_drift} passe à vrai. Ces résultats sont sauvegardés dans MLflow sous un run nommé \texttt{concept\_drift\_monitoring} et insérés dans la table ClickHouse \texttt{fact\_model\_drift} pour alerter l'équipe MLOps en vue d'un ré-entraînement.


\chapter{Orchestration, CI/CD et Déploiement}

\section{Orchestration sous Apache Airflow}
Les différents pipelines de données sont orchestrés de manière modulaire à l'aide de DAGs (Directed Acyclic Graphs) Airflow s'exécutant dans des conteneurs via \texttt{DockerOperator} :
\begin{enumerate}
    \item \textbf{DAG \texttt{spark\_mqtt\_to\_hudi}} : Gère le streaming Spark continu pour stocker les logs IoT bruts dans la table Hudi Bronze.
    \item \textbf{DAG \texttt{spark\_bronze\_to\_silver}} : Tâche batch s'exécutant toutes les 30 minutes. Il extrait les données Bronze, applique les baselines d'imputation et écrit la table Silver. Ce DAG déclare le dataset Silver en outlet.
    \item \textbf{DAG \texttt{spark\_silver\_to\_gold\_ml}} : Déclenché dynamiquement grâce au mécanisme \textbf{Data-Aware Scheduling} dès que le dataset Silver est mis à jour. Il coordonne l'inférence prédictive (Gold) et la détection subséquente de dérive statistique.
\end{enumerate}

\section{Intégration et Déploiement Continus (CI/CD)}
Afin de garantir la stabilité de la plateforme décisionnelle et l'absence de régression, un pipeline de CI/CD complet est automatisé via **GitHub Actions** :

\begin{figure}[H]
    \centering
    \begin{tikzpicture}[node distance=2cm, every node/.style={fill=green!10, draw, rounded corners, font=\footnotesize, align=center}]
        % Git push
        \node[fill=orange!20] (push) {\textbf{Git Push / PR}};
        
        % CI Steps
        \node[right=of push] (lint) {\textbf{Intégration Continue (CI)} \\ \hline Lintage : Flake8 \\ Formatage : Black \\ Validation : Airflow DAGs};
        
        % CD Steps
        \node[right=of lint, fill=blue!10] (build) {\textbf{Déploiement Continu (CD)} \\ \hline Build Docker Images \\ Publish to GHCR \\ (Registry)};
        
        % Target Deploy
        \node[right=of build, fill=red!10] (deploy) {\textbf{Production} \\ \hline docker-compose pull \\ Relancer les conteneurs};
        
        % Paths
        \draw[->, thick] (push) -- (lint);
        \draw[->, thick] (lint) -- (build);
        \draw[->, thick] (build) -- (deploy);
    \end{tikzpicture}
    \caption{Pipeline CI/CD avec GitHub Actions}
\end{figure}

Le pipeline se déroule en deux grandes phases :
\begin{itemize}
    \item \textbf{Phase de CI (Intégration Continue)} :
    \begin{itemize}
        \item **Lintage de syntaxe** : L'outil \texttt{Flake8} analyse le code Python des dossiers \texttt{spark/apps/}, \texttt{ml/} et \texttt{tools/} pour détecter les erreurs structurelles, de typage ou de non-respect de la norme PEP 8.
        \item **Formatage automatique** : \texttt{Black} s'assure de l'uniformité stylistique du code source.
        \item **Tests Airflow** : Un script de test unitaire charge l'ensemble des fichiers DAGs d'Airflow dans un interpréteur pour s'assurer de l'absence de cycles ou d'erreurs d'importation.
    \end{itemize}
    \item \textbf{Phase de CD (Déploiement Continu)} :
    \begin{itemize}
        \item **Compilation d'images Docker** : Si les tests CI réussissent, GitHub Actions compile de nouvelles versions des images Docker personnalisées. En particulier, l'image \texttt{iot-spark-hudi} (contenant Spark 3.5.6 et Apache Hudi 0.15.0) et l'image du bridge MQTT sont construites à l'aide des fichiers Dockerfile associés.
        \item **Publication GHCR** : Les images construites sont étiquetées (\textit{tagged}) et poussées sur le registre privé **GitHub Container Registry (GHCR)** sous l'adresse \texttt{ghcr.io/votre-depot/}.
        \item **Déploiement en Production** : Le serveur de production tire automatiquement les nouvelles images à l'aide de la commande \texttt{docker-compose pull} et relance les conteneurs impactés de manière transparente.
    \end{itemize}
\end{itemize}

\section{Déploiement de l'Infrastructure}
L'architecture complète du projet est déployée sous forme de services containerisés via Docker Compose :
\begin{itemize}
    \item \textbf{Base métadonnées Airflow} : Postgres 16.
    \item \textbf{MinIO \& MinIO MC} : Stockage et initialisation du bucket \texttt{iot-lake}.
    \item \textbf{Spark Master et 8 Workers} : Déploiement distribué totalisant 32 cœurs de calcul et 16 Go de mémoire RAM.
    \item \textbf{Mosquitto \& Bridge} : Broker de messagerie MQTT.
    \item \textbf{ClickHouse} : Entrepôt décisionnel pour le stockage Gold.
    \item \textbf{MLflow} : Serveur d'expériences.
\end{itemize}


\begin{center}
\addcontentsline{toc}{chapter}{Conclusion Générale}
\chapter*{Conclusion Générale}
\end{center}

Ce projet a permis de concevoir et de réaliser une infrastructure décisionnelle et MLOps robuste pour la maintenance prédictive de capteurs connectés industriels. En intégrant des formats de table transactionnels (Apache Hudi) à des serveurs OLAP haute performance (ClickHouse), le système garantit des temps de traitement optimisés ainsi qu'une fraîcheur élevée de l'information décisionnelle.

L'apport combiné du registre de modèles MLflow et de la détection statistique de dérive statistique par test de Kolmogorov-Smirnov assure une surveillance continue de la pertinence des prédictions de RUL, permettant de parer aux baisses de performance inhérentes aux changements d'environnement (dérive de données).

Comme perspectives, ce projet peut évoluer selon deux axes principaux :
\begin{enumerate}
    \item \textbf{Active Learning (Ré-entraînement automatique)} : Connecter l'alerte de dérive à un déclenchement automatisé du script d'entraînement pour adapter le modèle aux nouvelles données environnementales.
    \item \textbf{Modèles Temporels Profonds} : Remplacer l'algorithme XGBoost par des modèles d'apprentissage profond (LSTM, GRU ou Transformers) afin de mieux modéliser la dynamique temporelle sur de très longues séquences de vol.
\end{enumerate}

\end{document}
